\chapter{Fast-Koinzidenzkreis}

\section{Kontrolle der Fast Pulse}
Der Aufbau zur Kontrolle der Fast-Pulse ist in \fig{tikz/fast_pulse_kontrolle} dargestellt.

\jafps{tikz/fast_pulse_kontrolle}{Schaltplan zur Kontrolle der Fast Pulse}{0.75}

\subsection*{Linke Seite}
\jafps{links/fast_pulse}{}{0.9}
\todo{fast pulse vom scint kontrollieren}
\todo{polaritäten diskutieren}
\todo{anstiegszeit diskutieren}

\subsection*{Rechte Seite}


\section{CFD-Schwelle}
Der Aufbau, welcher zum Einstellen der CFD Schwellen genutzt wurde ist in \fig{tikz/cfd_schwelle_oszi} und der Aufbau zur Aufnahme eines Energiespektrums mit CFD als Gate in \fig{tikz/cfd_schwelle_mca} dargestellt.

\jafpp{tikz/cfd_schwelle_oszi}{}{tikz/cfd_schwelle_mca}{}
\subsection*{Linke Seite}
\todo{cfd schwelle}
\todo{erläutern was sich ändert wenn an cfd schwelle gedreht wird}
\jafps{links/scope_cfd_down}{Getriggert auf dem nicht gezeigten Kanal. Hier zu sehen die Signale, am DA, bei welchen der CFD getriggert hat.}{0.9}
\todo{ba spektrum mit cfd trigger}
\todo{cfd schwelle daraus?}
\todo{ist es sinnvoll ins spektrum zu schneiden}

\subsection*{Rechte Seite}


\section{Fast Koinzidenz}
Der Aufbau um die Fast-Kreis Koinzidenz zu einzustellen und zu überprüfen ist in \fig{tikz/fast_koinzidenz} abgebildet.
\jafps{tikz/fast_koinzidenz}{}{0.75}
\todo{fast koinzidenz}
\todo{AULFÖSUNGSZEIT vom TAC!!!}
\todo{welche quelle wird benutzt}
\todo{welcher detektor start/stop und warum}
\todo{relative zeitliche Lage von start und stop}
\todo{abstand von start zu stop}
\todo{wie lang und hoch sind die signale}

\section{Slow-Fast-Koinzidenz}
Der Aufbau um die Koinzidenz von Slow und Fast Kreis zu überprüfen und ggf. zu Korrigieren ist in \fig{tikz/fast_slow_koinzidenz} abgebildet.
\jafps{tikz/fast_slow_koinzidenz}{}{0.75}
\todo{fast-slow koinzidenz}
\todo{signale beschreiben}

\section{Promptkurve}
\todo{wie entsteht promptkurve}
\todo{warum natrium?}

